\input{preamble}

\begin{document}

\title{Stanley-Reisner schemes}
\maketitle

\phantomsection
\label{section-phantom}

\tableofcontents

\section{Stanley-Reisner schemes}
\label{section-sr}

\noindent
Here I follow \cite{luk1} for the basic definitions; these
are consistent with other sources.

Let $[n]:=\{0,\ldots,n\}$ be the set of all
positive integers from 0 to $n$, and let
$\Delta_n$ denote the power set of $[n]$. A
subset $\mathcal{K} \subseteq \Delta_n$ is
called a {\it simplicial complex} if for
every $f \in \mathcal{K}$ all the subsets of
$f$ are also in $\mathcal{K}$. Elements of
$\mathcal{K}$ are called {\it faces}.

Let $S=k[x_0,\ldots,x_n]$ be the polynomial ring in $n+1$ variables
over an algebraically closed field $k$.
If $a=\{i_1,\ldots,i_k\}\in \Delta_n$ we write $x_a \in S$ 
for the square-free monomial $x_{i_1}\cdot\ldots\cdot x_{i_k}$.
The {\it Stanley-Reisner} ideal
of a simplicial complex $\mathcal{K}$ is
$$
I_{\mathcal{K}}=\langle x_p : p \in \Delta_n\setminus \mathcal{K} \rangle
\subseteq S.
$$
The {\it Stanley-Reisner ring} of $\mathcal{K}$ is
$A_{\mathcal{K}}=S/I_{\mathcal{K}}$
and the {\it Stanley-Reisner scheme} is $\text{Proj} A_{\mathcal{K}}$.

\section{Stanley-Reisner scheme of Grünbaum sphere and a determinantal CY3}
\label{section-grunbaum-sphere}

\noindent
A Gröbner basis for the Stanley-Reisner ideal
of Grünbaum-Shreedharan's strange complex
$\mathcal{M}$ is 
$$
\begin{aligned}
x_6x_7x_8,
x_4x_6x_8, x_3x_7x_8, x_3x_5x_7, x_3x_4x_8,
x_2x_7x_8, \\ 
x_2x_5x_7, x_2x_5x_6,
x_2x_4x_7, x_2x_4x_6, x_1x_4x_6,\\ 
x_1x_4x_5,
x_1x_3x_8, x_1x_3x_6, x_1x_3x_5, x_1x_2x_5
\end{aligned}
$$
as computed in gfan and also GAP, and also as
shown in \cite{luk1}.

\section{GN variety}
\label{section-gn}

\noindent
There is a Calabi-Yau threefold in
$\mathbb{P}^7$ of degree 20 that could be the
smoothing of $\text{SR}(\mathcal{M})$. (As
far as I know there could exist others,
though they aren't very well-known.) It was
first studied by Gulliksen-Neg\aa ard
\cite{Gulliksen-Negard}. Among the papers I
read often for this project, it is mentioned
in \cite[Table 1 and last
paragraph]{CGKK-Aritmethically-Gorenstein-CY-in-P7}
and  \cite{cascade}. Additionally, this
variety has been studied in mirror-symmetry
context, see \cite{Jockers}.

In trying to immitate the computations I did
for the twisted cubic, I would like to find
the equations of the GN variety. In
\cite{cascade} we see its description as the
``$3\times 3$ minors of $4\times 4$ matrix
with linear forms''.

Back in May 19, 2025 I computed what might be ideal of the GN variety.
That is, as the $3 \times 3$ minors of the matrix
\begin{equation}
\label{equation-matrix}
\begin{bmatrix}
x_1 & x_2 & x_3 & x_4 \\
x_5 & x_6 & x_7 & x_8 \\
x_2 & x_3 & x_4 & x_5 \\
x_6 & x_7 & x_8 & x_1
\end{bmatrix}
\end{equation}
Those minors are
\[
\begin{aligned}
&x_3^2 x_5 - x_2 x_4 x_5 - x_2 x_3 x_6 + x_1 x_4 x_6 + x_2^2 x_7 - x_1 x_3 x_7, \\
&-x_3 x_6^2 + x_3 x_5 x_7 + x_2 x_6 x_7 - x_1 x_7^2 - x_2 x_5 x_8 + x_1 x_6 x_8, \\
&-x_3^2 x_6 + x_2 x_4 x_6 + x_2 x_3 x_7 - x_1 x_4 x_7 - x_2^2 x_8 + x_1 x_3 x_8, \\
&x_4 x_6^2 - x_4 x_5 x_7 - x_3 x_6 x_7 + x_2 x_7^2 + x_3 x_5 x_8 - x_2 x_6 x_8, \\
&x_3 x_4 x_5 - x_2 x_5^2 - x_2 x_4 x_6 + x_1 x_5 x_6 + x_2^2 x_8 - x_1 x_3 x_8, \\
&-x_1 x_2 x_5 + x_1^2 x_6 - x_4 x_6^2 + x_4 x_5 x_7 + x_2 x_6 x_8 - x_1 x_7 x_8, \\
&-x_1 x_2^2 + x_1^2 x_3 - x_3 x_4 x_6 + x_2 x_5 x_6 + x_2 x_4 x_7 - x_1 x_5 x_7, \\
&x_1 x_3 x_5 - x_1 x_2 x_6 + x_5 x_6^2 - x_5^2 x_7 - x_3 x_6 x_8 + x_2 x_7 x_8, \\
&x_4^2 x_5 - x_3 x_5^2 - x_2 x_4 x_7 + x_1 x_5 x_7 + x_2 x_3 x_8 - x_1 x_4 x_8, \\
&-x_1 x_3 x_5 + x_1^2 x_7 - x_4 x_6 x_7 + x_4 x_5 x_8 + x_3 x_6 x_8 - x_1 x_8^2, \\
&-x_1 x_2 x_3 + x_1^2 x_4 - x_4^2 x_6 + x_3 x_5 x_6 + x_2 x_4 x_8 - x_1 x_5 x_8, \\
&x_1 x_4 x_5 - x_1 x_2 x_7 + x_5 x_6 x_7 - x_5^2 x_8 - x_4 x_6 x_8 + x_2 x_8^2, \\
&x_4^2 x_6 - x_3 x_5 x_6 - x_3 x_4 x_7 + x_2 x_5 x_7 + x_3^2 x_8 - x_2 x_4 x_8, \\
&-x_1 x_3 x_6 + x_1 x_2 x_7 - x_4 x_7^2 + x_4 x_6 x_8 + x_3 x_7 x_8 - x_2 x_8^2, \\
&-x_1 x_3^2 + x_1 x_2 x_4 - x_4^2 x_7 + x_3 x_5 x_7 + x_3 x_4 x_8 - x_2 x_5 x_8, \\
&x_1 x_4 x_6 - x_1 x_3 x_7 + x_5 x_7^2 - x_5 x_6 x_8 - x_4 x_7 x_8 + x_3 x_8^2
\end{aligned}
\]
Now in theory I would just do all I did for the twisted cubic,
but now for GN, find a nice weight vector in for which the 
initial ideal is the ideal of $\text{SR}(\mathcal{M})$, and win.

But several problems arise:
\begin{enumerate}
\item First I have to double check that ideal for GN. I'm not
sure if those $3 \times 3$ minors of the $4 \times 4$ matrix I used are
good. {\bf Edit:} I think it's right. At least
by skimming through my references on GN, which are:
\cite{Jockers, Jockers2, Bertin, Gross-moduli-abelian}, I don't find
a matrix exactly but the claim is that a ``generic'' matrix
of $4 \times 4$ with linear forms in the entries should do the job.
(So that ChatGPT matrix above has indeed linear forms as entries.)
But most importantly \cite{Bertin} put a betti tally like
the ones one computes in Macaulay2 and indeed it's the same
as the one I get.

\item Most importantly, and this is what happened that other time
I tried to do this (when by the way I was being led by
ChatGPT to do exactly what I'm doing now, only now I
know a little more what's going on), is that the Gröbner fan
of GN is not computable. It's too much data and there's no output.
So it's not clear how to find the weight vector.
\end{enumerate}

One thing to try: obtain more freedom
on the $4\times 4$ matrix, see if you can change
that and still find the same good candidate.

\medskip\noindent
A breakthrough discovery:
the minimal resolution of GN and $\text{SR}(\mathcal{M})$ is not the same.

After this observation (and now knowing if
minimal resolutions of fibres of flat maps are isomorphic),
I'm not sure how to proceed. Possible ways are:

\begin{itemize}
\item Keep trying to prove/find out somehow 
(maybe asking: Eduardo Esteves thinks it's not true,
and recommended that I asked someone at URFJ)
whether of not the minimal resolution claim is true.
I may also look for other invariants on the fibers of a flat map,
see Stacks Project 
\href{https://stacks.math.columbia.edu/tag/02NM}{02NM}.

\item Keep trying to find the smoothing,
but I exhausted my computational attempts to find a perfect pairing
between the sets of Gröbner bases,
as well as my attempts to get around the size
of the Gröbner fan and just compute it via gfan.
\end{itemize}

\medskip\noindent
As a note, in \cite{luk2}, where the table
where GN was found lies, it is claimed that
there is  ``evidence'' that this table is
complete. The table counts arithmetically
Gorenstein CY3s. Adding up, either I find a
smoothing that is not the table, or prove
that it doesn't exist.


\section{6-vertex triangulation of RP2}
\label{section-6-vertex-triangulation-of-rp2}

\noindent
Here I read \cite{KuhnelBanchoff1983}.

\medskip\noindent
The icosahedron quotiented by antipodal
map gives a triangulation of 6 vertices
of $\mathbb{R}P^2$.
This can be embedded in $\mathbb{R}^5$ 
inside the 5-simplex $\Delta^5$,
which has 6 vertices.
This embedding has the properties
of symmetry (icosahedral$/-\text{Id}$),
duality (three vertices of $\Delta^5$ determine
a triangle of $\mathbb{R}P^2_6$ if and only if
the remaining do not),
tightness (too complicated) and
secant tangency (a chord joining
two points of $\mathbb{R}P^2$ is contained
in a boundary 4-simplex of $\Delta^5$.

Then consider Veronesse embedding
$[x_0:x_1:x_2]\mapsto [x_0^2:x_1^2:x_2^2:x_0x_1:x_0x_2:x_1x_2]$.
Like $\mathbb{R}P^2_6$ is contained in the boundary of $\Delta^5$,
this embedding is contained in $S^4$.
It also satisfies properties of symmetry (of $\mathbb{R}^3$),
duality, tightness and secant-tangency.

Each of the third and fourth properties
characterizes $\mathbb{R}P^2_v$ among
the smooth embeddings of $\mathbb{R}P^2$
in $\mathbb{R}^5$ not lying in an affine hyperplane.
``Pohl and Kuipar have in fact shown that
any {\it topological} embedding of the real
projective plane into $\mathbb{R}^5$ not lying in 
a hyperplane which has the tightness property
alone must either be $\mathbb{R}P^2_v$ or $\mathbb{R}P^2_v$.''

With respect to the complex projective plane,
the situation has been less satisfactory.

\section{9-vertex triangulation of CP2}
\label{section-9-vertex-triangulation-of-cp2}



\medskip\noindent
{\bf September 25.}
Summary of Calla's talk.
The objective is to construct a deformation 
of hyperkähler generalized Kummer fourfolds
starting from a special fiber.
The special fiber is the smoothing
of a Stanley-Reisner variety.

Proving this result would 
``imply that a general hyperkähler generalized kummer
fourfold with polarization of degree 2 and fixed divisibility
is contained in 36 quartics''.

\noindent
{\bf Questions.}
\begin{itemize}
\item How did they know that this funny simplicial
complex gives the special fiber of the degeneration
upon smoothing??
\item How does that result about being contained
in 36 quadric follows?
\item How come the special fiber of a degeneration
of fourfolds is a surface?
\end{itemize}

\medskip\noindent
Now let's go over the details of Calla's construction.
In \cite{KuhnelBanchoff1983} it was shown that there is a 9-vertex
triangulation $\Sigma$ of $\mathbb{C}P^2$.
This triangulation has 36 4-simplices.
Notice that the definition of SR scheme for Calla
might be different from ours: the {\it Stanley-Reisner variety}
of $\Sigma$ is the union of 36 copies of $\mathbb{P}^4$,
one for each 4-simplex of $\Sigma$.
Its ideal is generated by 36 quartics.

Part 1. We denote by $\Sigma$ the
9 vertex triangulation of  $\mathbb{P}^2$,
and $\text{SR}_{36}$ its corresponding
Stanley-Reisner scheme.

\begin{enumerate}

\item Blow up 9 coordinate points in $\mathbb{P}^8$ 
$(1:0:\ldots:0),\ldots,(0:\ldots:0:1)$.
(Recall that by the construction of the
Stanley-Reisner scheme associated to
a simplicial complex, in this case with 9
vertices, the vertices of the SR scheme
are precisely the coordinate points of $\mathbb{P}^8$.)
This introduces 9 copies of $\mathbb{P}^7$.
Proper transform of  $\text{SR}_{36}$
intersects $\mathbb{P}^7$ in a union
of  20 $\mathbb{P}^3$s, denoted $T_{20}$.
This does not depend on which copy of $\mathbb{P}^7$ 
we look at.

As a note, recall that the proper
transform of a subvariety $Y \subset X$ 
is defined as the {\it closure} of
$\pi^{-1}(Y\setminus Z)$ where
the blow up occurs at $Z$. In this case
$Z$ is the 9 vertices and $Y=\text{SR}_{36}$.
Calla claims that
that the closure of the proper
transform of $\text{SR}_{36}$ 
intersects each of the exceptional divisors
(one $\mathbb{P}^7$ for every vertex)
in a union of 20 $\mathbb{P}^3$ 
regardless of which exceptional divisor we
look at. We call that intersection $T_{20}$,
and it is inside $\mathbb{P}^7$.

\item Now we turn to modify $T_{20}$ 
to satisfy Friedman conditions.
Blow up 8 coordinate points in $\mathbb{P}^7$.
The proper transform of $T_{20}$ intersects
$\mathbb{P}^6$ in a union of 10 $\mathbb{P}^2$ s,
which we call $S_{10}$.
This {\bf does} depend on which copy of
$\mathbb{P}^7$ and $\mathbb{P}^6$ 
we are looking at!
$S_{10}$ is a Stanley-Reisner variety
and in fact is a triangulation of a sphere.
(A polyhedron!?)

\item Now we proceed to smooth $S_{10}$.
It has seven vertices, which we blow up,
introducing seven copies of $\mathbb{P}^5$.
Each of the $\mathbb{P}^2$ s becomes
a del Pezzo 6.
Looking at the picture we
see that there are vertices with
different valency. For each vertex
where $d$ $\mathbb{P}^2$s meet,
we glue in a del Pezzo surface of 
degree $d$. 
We get $\tilde{S}_{10}$, which is a 
simple normal crossing surface,
with trivial canonical bundle
and satisfies $d$-stability 
Friedman conditions.

Then there exists a Kulikov model
of type III degenerations of K3s
(of degree $10$ in $\mathbb{P}^7$)
with special fiber $\tilde{S}_{10}$.
(This is local smoothing over a
small disc, does not give a projective
family of K3s.)
\end{enumerate}

{\bf Key move.}
$S_{10}$ is a triangulation of a sphere.
Blow up the singularities and look
how does the proper transform of $S_{10}$ 
intersect the exceptional divisors.
To produce a smoothing,
we glue in del Pezzo surfaces
at the exceptional divisors.

Now we shall do the same for $T_{20}$.
(Which is still not the blow up
of $\text{SR}_{36}$, but the intersection
of the proper transform of $\text{SR}_{36}$ 
with any of the 9 exceptional divisors.)
We will glue in our projecive degeneration of K3s.
Notice that Friedman only gave us existence
of local model, so we must find equations.

For now let's just say that there's a bunch of procedures and
the resulting $\tilde{T}_{20}$ is a normal crossing and 
has trivial canonical bundle.
Then we can run smoothing procedure and find equations of family
of CY3s in $\mathbb{P}^7$ by noting that $T_{20} \subset 10 \mathbb{P}^4$s.
We should find generators of the ideal of
the variety $10\mathbb{P}^4$
among minors of $3\times 5$ matrix, similarly to surface case.

\medskip\noindent
{\bf October 19.}
It looks as if Calla was defining the SR variety of the triangulation
in a different way to my definition in Section \ref{section-sr}.
Today I checked that the ideal generated by 36 quartics
(one for every facet in the triangulation)
is not the same as the SR ideal.
How to tell which definition is ``correct''?
Correct for what? To answer I must get better understanding
of the deformation process from the Kummer perspective,
but it's not clear where this idea that the SR scheme
corresponding to the 9-vertex triangulation
would give the special fiber of a degeneration of
Kummer fourfolds!

I also noticed today that in \cite{KuhnelBanchoff1983}
there is a mention of $\mathcal{M}$!
I think $\mathcal{M}$ might be the vertex figure of $\mathbb{C}P^2_9$!


\section{Calla's talk}
\label{section-calla-talk}

\noindent
In this section I write literally Calla's talk.

\medskip\noindent
Aim: construct a type III
degeneration of varieties of generalised Kummer type 4-folds
with polarisation degree 2,
by starting from a given special fibre and constructing a smoothing.

Corollary. Proving this result will imply that a 
general HK variety of this type with polarisation
degree 2 (\& fixed divisibility) is contained
in 36 quartics.
(In particular, it is not contained in a cubic.)

\medskip\noindent
Reminder of what we mean by type III degen:

Figure with three strings of text:
\begin{itemize}
\item "Moduli space of gen. Kum. type
4-folds w/ fixed polarization" enclosed in a circle
in the center,
\item "Boundary of BB compactification
is made up of 1 dim components 
intersecting in dim 0" with an arrow
pointing at to lines intersecting the circle
of the last item, and
\item "Type III degen has special fibre corresponding
to this point", which has an arrow pointing
at the point of intersection of the last item.
\end{itemize}

\medskip\noindent
Method: 
\begin{itemize}
\item We give a candidate $\text{SR}_{36}$ for
the special fibre of this degen.
\item Since $\text{SR}_{36}$ is badly singular,
we make some modifications,
in order to be able to smoothe it.

\end{itemize}
\section{Moduli space of generalized Kummer type fourfolds}
\label{section-moduli-kummer}

\noindent
According to Calla's talk,
the whole point of constructing the smoothing
of $\text{SR}(\mathcal{M})$ is that it
should help us to smooth $\mathbb{C}P^2_9$,
which in turn should lie in the boundary
of the moduli space of generalized Kummer type fourfolds
with polarization degree 2.

Here's a quote from Calla's talk:
``Aim: construct a type III
degeneration of varieties of generalized
Kummer type 4-folds with polarization degree 2,
by starting from a given special fibre
and constructing a smoothing.''

\medskip\noindent
Let's address the following question.
Even if there exists a smoothing of $\text{SR}(\mathbb{C}P^2_9)$,
why would it be a generalized Kummer variety?

First we go over \cite[Theorem 0.2.1]{Fausk}:

\begin{theorem}
\label{theorem-fausk-021}
A smoothing, if it exists, of the Stanley-Reisner scheme
of a triangulation of the 3-sphere yields a Calabi-Yau 3-fold.
\end{theorem}

\begin{proof}
Two key steps: show using \cite[Theorem 6.1]{Bayer-Eisenbud}
that the canonical bundle of SR(sphere) is trivial.
Then show using Hartshorne that triviality of
canonical bundle stays the same under deformations.
\end{proof}

\noindent
Sure, that says that if I smooth $\text{SR}(M)$ 
then I have a CY 3fold. That's only the vertex-figure
of $\mathbb{C}P^2_9$.
Still, this looks far from saying that a smoothing
{\it of $\mathbb{C}P^2_9$} is a generalized Kummer 4fold.


\section{Dual complex, Hironaka theorem and Kummer type manifolds}
\label{section-dual-complex-hironaka-kummer-type-manifolds}

\noindent
These are notes from 
\href{http://verbit.ru/MATH/TALKS/Dual-complex-ddc_IMPA-2026.pdf}{
a seminar by Misha Verbitsky}
on his work with Nikon K.

\begin{definition}
\label{definition-snc}
Let $X \subset Y$ be a divisor in a smooth manifold.
We say that $X$ has normal crossings if
each of its components is smooth, and in a sufficiently
small neighbourhood of every point of $X$ the manifold
 $Y$ has coordinate system such that $X$ is identified with a union of
coordinate hyperplanes.
It has {\it simple normal crossings} if it has normal crossings and 
for each collection of irreducible components of the divisor,
their intersection is connected (possibly empty).
\end{definition}

\begin{theorem}[Hironaka resolution]
\label{theorem-hironaka}
Let $X$ be a complex variety.
Then there is a sequence of blow-ups with
smooth centers (=variety we are blowing up)
which results in a smooth resolution
$\pi:\tilde{X} \to X$ and the exceptional set of $\pi$ is SNC.
\end{theorem}

\begin{theorem}[Hironaka resolution of a pair]
\label{theorem-hironaka-pair}
Let $X$ be a complex variety and $Y \subset X$ a subvariety.
Then there is a sequence of blow-ups with smooth centers
which results in a smooth resultion
and the exceptional set of $\pi$ is a SNC divisor
and the preimage of $Y$ is also a SNC divisor.
\end{theorem}

\begin{definition}[Stepanov, 2004]
\label{definition-dual-complex}
Let $X \subset Y$ be a SNC divisor.
The {\it dual complex} of $X$ is a simplicial complex
whose vertices are the irreducible components $X_1,…X_n$,
and the $k$-simplices correspond to
intersections of $k+1$ of the irreducible components.
A {\it boundary} of a $k$-simplex is a $k-1$-simplex
obtained as the intersection of the divisors except one.
\end{definition}

\begin{theorem}[Stepanov]
Let $x \in X$ be an isolated singularity
and $\tilde{X}$ its smooth resolution such that
the preimage of $x$ is a SNC divisor.
Then the homotopy type of its dual complex is independent from
the choice of a resolution.
\end{theorem}

\noindent
Stepanov's theorem was generalized by
Arapura-Bakhtary-Wlodarczyk.

\begin{theorem}[ABW]
\label{theorem-abw}
Let $Y \subset X$ be a subvariety,
and $\tilde{X} \xrightarrow{\pi}X$ the resolution of the pair
$Y \subset X$.
Let $D$ be the exceptional set of $\pi$ (which is a SNC)
and $\tilde{Y}$ the set of all components
$D' \subset D$ such that $\pi(D')\subset Y$.
Then the homotopy type of the dual complex of $\tilde{Y}$
is independent from the choice of resoution.
\end{theorem}

\begin{definition}
\label{definition-kummer-type}
We say that $Y \subset X$ is of {\it Kummer type} 
if $b_1(C(X,Y))=0$ where $C(X,Y)$ is the dual complex
of this pair, and $b_1$ the rational Betti number.
\end{definition}

\begin{theorem}
\label{theorem-kummer-type}
Let $X \subset M$ be a complex subvariety of Kummer type,
and $G$ a finite group holomorphically acting on $M$
and preserving $X$.
Then $X/G \subset M/G$ is also Kummer type.
\end{theorem}

\medskip\noindent
Now we discuss $dd^c$-lemma and the dual complex.

\begin{theorem}
\label{theorem-kummer-type-manifolds}
Let $X \subset M$ be compact complex ($M$ is not necessarily compact).
Assume that $X$ is Fujiki class C
(i.e. bimeromorphic to Kähler) (normal) of Kummer type,
that is, $b_1(S)=0$ where $S$ is the dual complex
of the pair $X \subset M$.
If $\eta$ is an exact $(1,1)$-form,
then there exists a function $f$ defined
in a neighbourhood of $X$ such that $\eta=dd^cf$.
\end{theorem}

\begin{remark}
\label{remark-normality}
Normality of $X$ may not be used.
Recall that normal= when all meromorphic functions are bounded.
\end{remark}

\begin{proof}
\noindent
Step 1. Consider the resolution $\tilde{X} \subset \tilde{M}$
With some extra blow ups we may assume that all the 
components of $\tilde{X}$ are Kähler.
Then we apply $dd^c$-lemma on any of the components.

\medskip\noindent
Step 2. $dd^cf=0$ means that it is a sum of holomorphic
and antiholomorphic,
which means that $\ker dd^c$ are
constant functions. There's also maximum princple
which implies that $f$ restricted to $\pi^{-1}(x)$ for any
$x$ in $X$ is constant.

\medskip\noindent
Step 3. Choose a function $f_i$ which satisfies
$dd^c f_i=\eta|_{D_i}$ on each of these components.
Then we prove that this satisfies the cocycle condition,
and we thus from $\eta$ we obtain a canonical element 
$\mu_\eta \in H^1(S,\mathbb{C})$.
[I think this cocycle condition is used
to glue the $f_i$ to a global function].

\medskip\noindent
Step 4. Extending $f$ to a neighbourhood…
\end{proof}

\noindent
Then there is a statement about the $dd^c$ lemma not
holding if $X$ is not of Kummer type: there is a form
that is not $dd^c$-exact:

\begin{theorem}
\label{theorem-ddc-counterexample}
Let $X \subset M$be a compact, positive-dimensional
subvariety of a manifold $M$, which is not of Kummer type.
Assume that there exists a non-contractible path in the
dual complex of the resolution of the pair, $S$,
which visits a component $D_0 \subset \tilde{X}$
only once, and the projection of $D_0$ to $M$
is positive-dimensional.
Assume $\dim M=2$.
Then there exists an exact $(1,1)$-form $\eta$
on a neighbourhood of $X$ such that $\eta \neq  dd^cf$ for any 
function $f$ on $X$.
\end{theorem}

\medskip\noindent
For the proof of ABW theorem we need the following results,
which were proved in 2002 in D. Abramovish, K. Karu, K. Matsuky,
J. Wlodarczyk, Torication and factorization of birational maps.

\begin{theorem}[Weak factorization theorem]
\label{theorem-weak-factorization-theorem}
Let $\varphi: X \dashrightarrow X'$ be a bimeromorphic map
of compact complex manifolds and $U \subset X$, $U' \subset X'$
the complements to the exceptional sets.
Then $\varphi$ can be decomposed 
into a sequence of blow-ups and blow-downs with
smooth centers.
\end{theorem}

\begin{remark}
\label{remark-meromorphic}
So what does it mean to be meromorphic? It means that there is
$$
\xymatrix{
&  Z \subset X \times X'\ar[dl]\ar[dr]\\
X & & X'
}
$$
where projections are proper, dominant and 1-to-1. [Maybe something
else I couldn't write down.
\end{remark}

\noindent
We need a version with boundary:

\begin{theorem}
\label{theorem-weak-factorization-for-pairs}
Let $\varphi:X \dashrightarrow X'$ be a bimeromorphic map
of compact complex manifolds, and $U \subset X$, $U' \subset X'$
be complements to the exceptional sets.
Consider SNC divisors $D \subset X$, $D' \subset X'$
and let $D_1,…, D_n$ and $D_1',…,D_n'$
be their irreducible components.
Assume that $\varphi$ takes $X\setminus D_i$
to $X' \setminus D_i'$
and this bimeromorphic map is proper
(that is, the closure of its grap in 
$X \setminus D_i \times X'\setminus D'$
projects properly to the components $X \setminus D_i$
and $X'\setminus D_i$).
Then $\varphi$ can be decomposed into a sequence
of blow-ups and blow-downs with smooth
centers, and the exceptional sets of blow-ups
have SNC with the proper preimages of $\cup  D_i$.
\end{theorem}

\noindent
Now we can prove Theorem \ref{theorem-abw},
which says that the homotopy type of the dual complex
is independent from the resolution.

\begin{proof}[Proof of Theorem \ref{theorem-abw}]
\noindent
Step 1. Take to resolutions of the pair.
Then they are bimeromorphic.
If we can show that the homotopy type of the dual
complex stays the same in each of the steps of the 
weak factorization theorem, then we are done.

\medskip\noindent
Step 2. We analyse the different possible cases of center in the 
first step. The first case is easy: the dual complex remains the same.
The second case I didn't understand.
The third case boils down to gluing a simplex to the 
original dual complex.
\end{proof}

\noindent
Then there is a $G$-equivariant version of this theorem
and of the weak factorization theorem.

\medskip\noindent
Finally we sketch Theorem \ref{theorem-kummer-type}.
Using $G$-invariant Hironaka, we
resolve, take quotient, take product, resolve, etc…
This would give a sequence of dual complexes…

\medskip\noindent
Take Hilbert scheme of torus, projected to symmetric
product of elliptic curve…
$$
\xymatrix{
(T^n)^{[n]}\ar[d]\\ \text{Sym}E.
}
$$


\section{Dual complex of a degeneration}
\label{section-dual-complex}

\noindent 
Here I review \cite{KLSV} to
understand the relation between
$\mathbb{C}P_9^2$ and the moduli of
generalized Kummer varieties.

``It turns out that the right class of
singularities that still allow the definition
of a meaningful dual complex is dlt. In other
words, the correct context for defining and
intrincis dual complex associated to a
degeneration is that of minimal dlt
degeneration. The minimal dlt model
$\mathcal{X}/\Delta$ is not unique, but
changing the model has no effect on
$|\Sigma|$.''

``Relevant for us is the fact that associated
to a $K$-trivial [I guees, trivial canonical
bundle] degenertion $\mathcal{X}/\Delta$
there is an intrinsic (depending only on
$\mathcal{X}^* /\Delta^*$) topological space
that we call (following [MN15]) {\it the
essential skeleton} $\text{Sk}(X)$ associated
to the degeneration. For a minimal dlf
degeneration of $K$-trivial varieties, the
essential skeleton $\text{Sk}(\mathcal{X})$
can be identified with the topological
realization $|\Sigma|$ of the dual complex
(cf Nicaise-Xu).''

``The purpose of this section is to make some
remarks on the structure of the essential
skeleton $\text{Sk}(\mathcal{X})$ for a
degeneration of hyper-Kähler manifolds''.

``[…] we note that the
$\text{Sk}(\mathcal{X})$ for a MUM
degeneration of Calaby-Yau $n$-folds is
predicted to be the sphere $S^n$. This is
true in dimension 2 by Kulikov's theorem, and
in dimensions 3 (unconditional) and 4
(assuming additionally that degeneration is
normal crossings by Kollár-Xu.''

``6.2.6. To conclude, mirror symmetry (via
SYZ conjecture and Kontsevich-Soibelman)
predicts that the essential skeleton
$\text{Sk}(\mathcal{X})$ for a MUM
degeneration is $S^n$ and respectively
$\mathbb{C}P^n$ for Calabi-Yau's and
respectively hyper-Kähler's. The following
result is a weaker version of this statement,
saying that it holds in a cohomological
sense.''

\begin{theorem}
\label{theorem-skeleton-cohomology}
\begin{reference}
\cite[Theorem 6.14]{KLSV}
\end{reference}
Let $\mathcal{X}/\Delta$ be a minimal dlt
degeneration of $K$-trivial varieties.
Assume that the general fiber $X_t$ 
is a simply connected $K$-trivial variety.

Assume that the degeneration is maximal unipotent.
Then
\begin{enumerate}
\item $H^* (\text{Sk}(\mathcal{X}),\mathbb{Q})
\cong \prod_iH^* (S^{k_i},\mathbb{Q})
\times_jH^* (\mathbb{C}P^{l_j},\mathbb{Q}))$,
where $k_i$ represent the dimensions
of the Calabi-Yau factos and $2l_j$ 
the dimensions of the hyper-Kähler factos
in the Beauville-Bogomolov decomposition of the general fiber $X_t$.

\item Conversely, the cohomology algebra
of the skeleton $\text{Sk}(\mathcal{X})$ 
determines the type of the Beauville-Bogomolov
decomposition of $X_t$.
\end{enumerate}
\end{theorem}



\section{Degenerations of Hilbert schemes (notes on Calla's thesis)}
\label{section-degenerations-hilbert}

\noindent
Notes on Calla's thesis on
degenerations of Hilbert schemes after
\href{https://www.youtube.com/watch?v=Lf4NYTvnO9I&t=969s}
{this video}.

The aim is to construct ``good'' degenerations
of Hilbert schemes of points over semistable
degenerations of surfaces. This was motivated by
the study of the Hilbert schemes of points
on ``maximally degenerate K3 surfaces'',
i.e. K3 surfaces that are on the boundary of
the compactified moduli space of K3s.

Basic setup: a flat projective family
of surfaces $X \to C \cong \mathbb{A}^1$ over
an algebraically closed field of characteristic 0.
Consider Étale local model given by
$\Spec k[x,y,z,t]/(xyz-t)$.
The general fibers are smooth, and
the special fibre $X_0$ over $0 \in C$ 
has simple normal crossing singularity.
We picture this as the three axes $x,y,z$ in $X_0$ 
intersecting at the origin.

Let $X^0= X\setminus X_0$ and $C^0=C\setminus \{0\}$.
Consider the Hilbert scheme of each fiber.
We have removed the limit over 0.
The question is: how can we compactify the Hilbert scheme
$\text{Hilb}^n(X^0/C^0)$ after we removed that fiber?
How to add a ``good limit''?

Starting from a ``big'' deformation space,
namely a fiber product $X \times_{\mathbb{A}^1}\mathbb{A}^n$,
Calla produces $X[n]$ after a sequence of blow-ups,
and later takes a quotient using a group action.
Thus produces a stack $\mathfrak{X}$
that contains the modifications
on the singular fiber
as well as the original fibers.

\begin{theorem}[Calla]
\label{theorem-calla}
The stack of Li-Wu stable length in 0-dimensional
subschemes in $\mathfrak{X}$ is Deligne-Mumford
and proper over $C$.
\end{theorem}

\noindent
It turns out this is a specific example
where things work out better than expected,
essentially by the way $X[n]$ was constructed.
The general case is not so ``nice''
(see [Maulik-Ranganathan]).

\medskip\noindent
{\bf Note.} At some point the main focus
of Calla's work turned to the Hilbert schemes.
I wonder if the generalised Kummer fourfolds studied in the
$\mathbb{C}P^2_9$ topic are actually examples
of these Hilbert schemes, or if it is a genuine
change of subject.


\end{document}
